% ==============================================================================
% Nature journal submission template (skeleton)
% Nature does not ship an official document class on CTAN; the production
% template is distributed via Nature's author services and Overleaf. This
% skeleton uses the standard `article` class with Nature-editorial styling so
% it compiles on any TeX Live install. Before final submission, download the
% official Nature LaTeX template from Overleaf ("Nature template for papers")
% or from https://www.nature.com/nature/for-authors/final-submission and
% replace the preamble below.
%
% Editorial style (Nature, "Author submissions" guidelines):
%   Layout     : single column, double-spaced, line numbers for review
%   Body font  : Times Roman 11-12pt (Arial/Helvetica acceptable)
%   Math font  : Times-compatible (mathptmx)
%   Captions   : 8-9pt, sentence case
%   Citation   : superscript numerals in text, numbered reference list
%   Figures    : single column 89 mm  |  double column 180 mm
%   Sections   : Introduction (no heading), Results, Discussion, Methods
%   Abstract   : one paragraph, ~150 words, no heading, bold lead
% References  : nature.bst (numeric, superscript), via natbib `super` option
% ==============================================================================
\documentclass[11pt,a4paper]{article}

% ---- Page geometry: A4 with generous review margins ----
\usepackage[a4paper,top=25mm,bottom=25mm,left=25mm,right=25mm]{geometry}

% ---- Fonts: Times text + Times math (Nature preference) ----
\usepackage[T1]{fontenc}
\usepackage[utf8]{inputenc}
\usepackage{mathptmx}             % Times for body AND math
\usepackage{helvet}               % Helvetica/Arial as sans companion
\usepackage{courier}

% ---- Line numbers + double spacing for review ----
\usepackage{lineno}               % \linenumbers  for review mode
\usepackage{setspace}
\doublespacing                    % comment out for camera-ready

% ---- Math, figures, tables ----
\usepackage{amsmath,amssymb}
\usepackage{graphicx}
\usepackage{booktabs}
\usepackage{array}
\usepackage{xcolor}

% ---- Citations: Nature superscript numerals ----
\usepackage[super,sort&compress,comma]{natbib}
\bibliographystyle{unsrtnat}      % numeric, citation-order; switch to nature.bst if available
% If you have nature.bst on your system, replace the line above with:
%   \bibliographystyle{nature}

% ---- Hyperref last ----
% \usepackage{hyperref}
% \hypersetup{colorlinks=true,linkcolor=blue,citecolor=blue,urlcolor=blue}

% ---- Nature-style section formatting (no section numbering) ----
\setcounter{secnumdepth}{-1}      % unnumbered section headings; use \section*{} in body

\title{\textbf{A concise, declarative title of the Nature submission}}
\author{First Author\textsuperscript{1,*}, Second Author\textsuperscript{2}, Senior Author\textsuperscript{1,*}\\
\textsuperscript{1}Department of X, University of Y, City, Country.\\
\textsuperscript{2}Department of Z, Institute of W, City, Country.\\
\textsuperscript{*}e-mail: first.author@univ.edu; senior.author@univ.edu}
\date{}

\begin{document}
\linenumbers                     % review mode: show line numbers

\maketitle

% ==============================================================================
% Abstract (Nature: one paragraph, ~150 words, no "Abstract" heading, bold lead)
% ==============================================================================
\noindent\textbf{Abstract.} % TODO: one paragraph, ~150 words. Lead with the
% headline finding in bold, then the question, the approach, the key result
% with effect sizes, and the broader implication. No citations.

% ==============================================================================
% Introduction (Nature: no section heading, opens with the broad question)
% ==============================================================================
\section*{Introduction}
% TODO: 3-5 paragraphs. (i) broad context and why it matters, (ii) specific
% gap in current knowledge, (iii) conceptual approach of this work, (iv) the
% main finding stated up-front (Nature style: results preview).

% ==============================================================================
% Results (Nature: numbered sub-headings, figures referenced inline)
% ==============================================================================
\section*{Results}
% TODO: open with a one-paragraph summary of the main finding; then sub-sections
% each beginning with a topic sentence that states the result before the data.

\subsection*{Result one: a short, declarative sub-heading}
% TODO: present the result, reference the figure inline, e.g. (Fig. 1).
% Describe controls and statistical support (n, test, p-value, effect size).

\subsection*{Result two: another declarative sub-heading}
% TODO: second result cluster.

% ==============================================================================
% Discussion (Nature: 2-4 paragraphs, no numbered sub-headings)
% ==============================================================================
\section*{Discussion}
% TODO: (i) restate the main result and how it advances the field,
% (ii) compare with prior work and reconcile differences, (iii) limitations,
% (iv) future directions and broader implications.

% ==============================================================================
% Methods (Nature: detailed, reproducible; can be smaller type after refs)
% ==============================================================================
\section*{Methods}
% TODO: data sources, experimental design, statistical analysis, code/data
% availability. Nature requires enough detail for replication.

% ==============================================================================
% References  —  numbered, superscript in text
% Compile:  pdflatex main  ->  bibtex main  ->  pdflatex main  ->  pdflatex main
% ==============================================================================
\bibliography{references}
% Inline fallback (uncomment if not using BibTeX):
% \begin{thebibliography}{99}
% \bibitem{ref_example} Author, F. \& Author, S. Title of the paper.
%   \emph{Nature} \textbf{XXX}, 1--10 (20YY).
% \end{thebibliography}

% ==============================================================================
% Acknowledgements, Author Contributions, Competing Interests, Data Availability
% ==============================================================================
\section*{Acknowledgements}
% TODO: funding agencies, grant numbers, individuals.

\section*{Author Contributions}
% TODO: e.g. "F.A. designed the study, performed the analysis, and wrote the
% paper. S.A. collected the data. S.A. supervised the project."

\section*{Competing Interests}
% TODO: declare any competing interests, or state "The authors declare no
% competing interests."

\section*{Data \& Code Availability}
% TODO: state where data and code are deposited (with DOIs / repository links).

% ==============================================================================
% Figure Legends (Nature: legends grouped after the main text)
% ==============================================================================
\section*{Figure Legends}

\noindent\textbf{Figure 1.} % TODO: title in bold, then a self-contained
% description of what is plotted, the sample size, statistical test, and
% error bars (e.g. "mean $\pm$ s.d., $n=10$ biological replicates").

\end{document}
